\documentclass[12pt,a4paper]{article}
\usepackage[a4paper,top=15mm,bottom=15mm,left=16mm,right=16mm]{geometry}
\usepackage[T1]{fontenc}
\usepackage{amsmath}
\usepackage{amssymb}
\usepackage{enumitem}
\usepackage{xcolor}

\setlength{\parindent}{0pt}
\setlength{\parskip}{4pt}
\allowdisplaybreaks

\newcommand{\ptitle}[1]{%
  \par\addvspace{12pt}
  \noindent{\color{blue!50!black}\rule{\linewidth}{0.8pt}}\par\nobreak\vspace{2pt}
  \noindent{\large\bfseries\color{blue!50!black}#1}\par\nobreak\vspace{2pt}
  \noindent{\color{blue!50!black}\rule{\linewidth}{0.4pt}}\par\vspace{8pt}}

\newcommand{\wsol}{%
  \par\smallskip
  \noindent\textbf{\color{blue!50!black}Worked solution.}\space}

\begin{document}

\begin{center}
{\Large\bfseries Recurring Decimals into Fractions}\\[3pt]
{\large Worked Solutions}
\end{center}

\ptitle{Part 1: Reading and Writing Recurring Decimals}

\noindent\textbf{Q1.} Write out the first eight decimal places.
\wsol
\begin{enumerate}[label=(\alph*),leftmargin=*]
\item $0.\dot{5}=0.55555555$
\item $0.\dot{1}\dot{8}=0.18181818$
\item $0.3\dot{7}=0.37777777$
\item $0.\dot{2}0\dot{6}=0.20620620$
\item $0.4\dot{5}\dot{1}=0.45151515$
\item $0.\dot{9}=0.99999999$
\end{enumerate}

\medskip
\noindent\textbf{Q2.} Write each using dot notation.
\wsol
\begin{enumerate}[label=(\alph*),leftmargin=*]
\item $0.8888\ldots=0.\dot{8}$
\item $0.636363\ldots=0.\dot{6}\dot{3}$
\item $0.24444\ldots=0.2\dot{4}$
\item $0.157157\ldots=0.\dot{1}5\dot{7}$
\item $0.91666\ldots=0.91\dot{6}$
\item $0.40505\ldots=0.4\dot{0}\dot{5}$
\end{enumerate}

\medskip
\noindent\textbf{Q3.}
\wsol
\begin{enumerate}[label=(\alph*),leftmargin=*]
\item We have
\[
0.\dot{4}\dot{5}=0.454545\ldots
\qquad\text{and}\qquad
0.4\dot{5}=0.455555\ldots
\]
They agree in the first two decimal places. At the third decimal place, the first number has $4$ and the second has $5$. Hence $0.4\dot{5}$ is bigger.

\item There are many possible answers. For example,
\[
0.\dot{3}\dot{5}=0.353535\ldots
\]
is bigger than $0.\dot{3}=0.3333\ldots$ and smaller than $0.\dot{4}=0.4444\ldots$.
\end{enumerate}

\ptitle{Part 2: Finding Fractions from Recurring Decimals}

\noindent\textbf{Q4.} Fill in the gaps to turn $0.\dot{7}$ into a fraction.
\wsol
\textbf{Step 1:} Let $x=0.\dot{7}=0.7777\ldots$.

\textbf{Step 2:} Multiplying by $10$ shifts every digit one place to the left:
\[
10x=7.7777\ldots
\]

\textbf{Step 3:} The tails are identical, so they cancel on subtraction:
\[
10x-x=7.7777\ldots-0.7777\ldots=7.
\]
Therefore
\[
9x=7.
\]

\textbf{Step 4:} Divide both sides by $9$:
\[
x=\frac{7}{9}.
\]

\medskip
\noindent\textbf{Q5.} Turn $0.\dot{8}$ into a fraction.
\wsol
Let $x=0.\dot{8}=0.8888\ldots$.
\[
10x=8.8888\ldots
\]
Subtract:
\[
10x-x=8,
\]
so
\[
9x=8.
\]
Hence
\[
x=\frac{8}{9}.
\]

\medskip
\noindent\textbf{Q6.} Turn each into a fraction in its simplest form.
\wsol
\begin{enumerate}[label=(\alph*),leftmargin=*]
\item $0.\dot{2}$: let $x=0.2222\ldots$. Then $10x=2.2222\ldots$, so $9x=2$ and $x=\frac{2}{9}$.

\item $0.\dot{5}$: let $x=0.5555\ldots$. Then $10x=5.5555\ldots$, so $9x=5$ and $x=\frac{5}{9}$.

\item $0.\dot{6}$: let $x=0.6666\ldots$. Then $10x=6.6666\ldots$, so $9x=6$ and $x=\frac{6}{9}=\frac{2}{3}$.

\item $0.\dot{3}$: let $x=0.3333\ldots$. Then $10x=3.3333\ldots$, so $9x=3$ and $x=\frac{3}{9}=\frac{1}{3}$.
\end{enumerate}

\ptitle{Part 3: Longer Repeating Blocks}

\noindent\textbf{Q7.} Fill in the gaps to turn $0.\dot{7}\dot{2}$ into a fraction.
\wsol
\textbf{Step 1:} Let $x=0.\dot{7}\dot{2}=0.727272\ldots$.

\textbf{Step 2:} The repeating block has $2$ digits, so multiply by $100$:
\[
100x=72.727272\ldots
\]

\textbf{Step 3:} The tails now match, so
\[
100x-x=72.727272\ldots-0.727272\ldots=72.
\]
Thus
\[
99x=72.
\]

\textbf{Step 4:} Divide both sides by $99$:
\[
x=\frac{72}{99}=\frac{8}{11}.
\]

\medskip
\noindent\textbf{Q8.} Turn these into fractions in their simplest form.
\wsol
\begin{enumerate}[label=(\alph*),leftmargin=*]
\item $0.\dot{1}\dot{5}$: let $x=0.151515\ldots$. Then $100x=15.151515\ldots$, so $99x=15$ and
\[
x=\frac{15}{99}=\frac{5}{33}.
\]

\item $0.\dot{0}\dot{6}$: let $x=0.060606\ldots$. Then $100x=6.060606\ldots$, so $99x=6$ and
\[
x=\frac{6}{99}=\frac{2}{33}.
\]

\item $0.\dot{5}\dot{4}$: let $x=0.545454\ldots$. Then $100x=54.545454\ldots$, so $99x=54$ and
\[
x=\frac{54}{99}=\frac{6}{11}.
\]

\item $0.\dot{9}\dot{0}$: let $x=0.909090\ldots$. Then $100x=90.909090\ldots$, so $99x=90$ and
\[
x=\frac{90}{99}=\frac{10}{11}.
\]

\item $0.\dot{1}0\dot{8}$: the repeating block is $108$, which has $3$ digits, so multiply by $1000$.
Let $x=0.108108\ldots$. Then $1000x=108.108108\ldots$, so $999x=108$ and
\[
x=\frac{108}{999}=\frac{4}{37}.
\]
\end{enumerate}

\medskip
\noindent\textbf{Q9.}
\begin{enumerate}[label=(\alph*),leftmargin=*]
\item To decide what to multiply by, count the number of digits in the repeating block. If the block has $n$ digits, multiply by $10^n$. This moves the decimal point $n$ places, so the repeating block lines up exactly with the block in the original decimal. The tails then cancel when the two decimals are subtracted.

\item $0.\dot{1}\dot{2}=0.121212\ldots$. Let $x=0.\dot{1}\dot{2}$. Then $100x=12.121212\ldots$, so
\[
99x=12
\qquad\text{and}\qquad
x=\frac{12}{99}=\frac{4}{33}.
\]
Ben has simply simplified Kayla's fraction. Both are correct.
\end{enumerate}

\ptitle{Part 4: When the Repeating Part Does Not Start Straight Away}

\noindent\textbf{Q10.} Fill in the gaps to turn $0.2\dot{7}$ into a fraction.
\wsol
Let $x=0.2\dot{7}=0.27777\ldots$.

Shift past the $2$:
\[
10x=2.7777\ldots
\]
Shift one more place:
\[
100x=27.7777\ldots
\]
Subtract the first from the second:
\[
100x-10x=27.7777\ldots-2.7777\ldots=25.
\]
Thus
\[
90x=25
\]
and
\[
x=\frac{25}{90}=\frac{5}{18}.
\]

\medskip
\noindent\textbf{Q11.} Turn these into fractions in their simplest form.
\wsol
\begin{enumerate}[label=(\alph*),leftmargin=*]
\item $0.4\dot{3}=0.4333\ldots$. Let $x=0.4333\ldots$.
\[
10x=4.3333\ldots,\qquad 100x=43.3333\ldots
\]
Subtract:
\[
90x=39
\]
so
\[
x=\frac{39}{90}=\frac{13}{30}.
\]

\item $0.58\dot{3}=0.58333\ldots$. Let $x=0.58333\ldots$.
\[
100x=58.3333\ldots,\qquad 1000x=583.3333\ldots
\]
Subtract:
\[
900x=525
\]
so
\[
x=\frac{525}{900}=\frac{7}{12}.
\]

\item $0.08\dot{3}=0.08333\ldots$. Let $x=0.08333\ldots$.
\[
100x=8.3333\ldots,\qquad 1000x=83.3333\ldots
\]
Subtract:
\[
900x=75
\]
so
\[
x=\frac{75}{900}=\frac{1}{12}.
\]

\item $0.1\dot{2}\dot{4}=0.1242424\ldots$. Let $x=0.1242424\ldots$.
\[
10x=1.242424\ldots,\qquad 1000x=124.242424\ldots
\]
Subtract:
\[
990x=124-1=123
\]
so
\[
x=\frac{123}{990}=\frac{41}{330}.
\]
\end{enumerate}

\medskip
\noindent\textbf{Q12.} Turn these into improper fractions in their simplest form.
\wsol
\begin{enumerate}[label=(\alph*),leftmargin=*]
\item $2.\dot{7}=2.7777\ldots$. Let $x=2.7777\ldots$.
\[
10x=27.7777\ldots
\]
Subtract:
\[
9x=25
\]
so
\[
x=\frac{25}{9}.
\]

\item $1.\dot{4}\dot{5}=1.454545\ldots$. Let $x=1.454545\ldots$.
\[
100x=145.454545\ldots
\]
Subtract:
\[
99x=144
\]
so
\[
x=\frac{144}{99}=\frac{16}{11}.
\]

\item $3.\dot{1}\dot{8}=3.181818\ldots$. Let $x=3.181818\ldots$.
\[
100x=318.181818\ldots
\]
Subtract:
\[
99x=315
\]
so
\[
x=\frac{315}{99}=\frac{35}{11}.
\]
\end{enumerate}

\medskip
\noindent\textbf{Q13.} Mixed practice.
\wsol
\begin{enumerate}[label=(\alph*),leftmargin=*]
\item $0.\dot{8}$:
$x=0.8888\ldots$, so $10x=8.8888\ldots$. Therefore $9x=8$ and $x=\frac{8}{9}$.

\item $0.\dot{6}\dot{3}$:
$x=0.636363\ldots$, so $100x=63.636363\ldots$. Therefore $99x=63$ and $x=\frac{63}{99}=\frac{7}{11}$.

\item $0.2\dot{5}$:
$x=0.2555\ldots$, so $10x=2.5555\ldots$ and $100x=25.5555\ldots$. Subtracting gives $90x=23$, so $x=\frac{23}{90}$.

\item $0.\dot{0}2\dot{7}$:
$x=0.027027\ldots$, so $1000x=27.027027\ldots$. Therefore $999x=27$ and $x=\frac{27}{999}=\frac{1}{37}$.

\item $0.7\dot{2}$:
$x=0.7222\ldots$, so $10x=7.2222\ldots$ and $100x=72.2222\ldots$. Subtracting gives $90x=65$, so $x=\frac{65}{90}=\frac{13}{18}$.

\item $0.\dot{1}\dot{2}$:
$x=0.121212\ldots$, so $100x=12.121212\ldots$. Therefore $99x=12$ and $x=\frac{12}{99}=\frac{4}{33}$.

\item $0.41\dot{6}$:
$x=0.41666\ldots$, so $100x=41.6666\ldots$ and $1000x=416.6666\ldots$. Subtracting gives $900x=375$, so $x=\frac{375}{900}=\frac{5}{12}$.

\item $0.06\dot{4}$:
$x=0.06444\ldots$, so $100x=6.4444\ldots$ and $1000x=64.4444\ldots$. Subtracting gives $900x=58$, so $x=\frac{58}{900}=\frac{29}{450}$.
\end{enumerate}

\ptitle{Part 5: Challenges}

\noindent\textbf{Q14.}
\begin{enumerate}[label=(\alph*),leftmargin=*]
\item Let $x=0.\dot{9}=0.9999\ldots$.
\[
10x=9.9999\ldots
\]
Subtract:
\[
10x-x=9
\]
so
\[
9x=9
\]
and hence
\[
x=1.
\]

\item A good reply to Ella is:

The algebra is not a trick. If $0.\dot{9}=x$, then $10x=9.\dot{9}$, so subtracting gives $9x=9$, and therefore $x=1$.

Also, $\frac13=0.\dot{3}$. Multiplying both sides by $3$ gives
\[
1=0.\dot{9}.
\]
On the number line, the numbers $0.9$, $0.99$, $0.999$, \ldots get closer and closer to $1$. Each new one closes half of the remaining gap. Since there is no positive gap left, $0.\dot{9}$ must actually equal $1$.

\item There is no largest number smaller than $1$. If someone names a number $a<1$, then
\[
a<\frac{a+1}{2}<1.
\]
So there is always another number between $a$ and $1$, which is bigger than $a$ but still smaller than $1$. This is also why $0.\dot{9}$ cannot be ``just below'' $1$: there is no room for a largest number below $1$.
\end{enumerate}

\medskip
\noindent\textbf{Q15.}
\wsol
\begin{enumerate}[label=(\alph*),leftmargin=*]
\item Using short division:
\[
\frac12=0.5,\qquad
\frac13=0.\dot{3},\qquad
\frac14=0.25,\qquad
\frac15=0.2,
\]
\[
\frac16=0.1\dot{6},\qquad
\frac17=0.\dot{1}4285\dot{7},\qquad
\frac18=0.125,
\]
\[
\frac19=0.\dot{1},\qquad
\frac1{10}=0.1,\qquad
\frac1{11}=0.\dot{0}\dot{9},\qquad
\frac1{12}=0.08\dot{3},\qquad
\frac1{16}=0.0625.
\]
The terminating decimals are
\[
\frac12,\quad\frac14,\quad\frac15,\quad\frac18,\quad\frac1{10},\quad\frac1{16}.
\]

\item The terminating denominators are:
\[
2=2,\quad 4=2^2,\quad 5=5,\quad 8=2^3,\quad 10=2\times5,\quad 16=2^4.
\]
They all have only the prime factors $2$ and $5$.

\item
\[
\frac1{15}\text{ recurs, since }15=3\times5.
\]
\[
\frac1{40}\text{ terminates, since }40=2^3\times5.
\]
\[
\frac1{35}\text{ recurs, since }35=5\times7.
\]
\[
\frac1{64}\text{ terminates, since }64=2^6.
\]

\item A terminating decimal can be written with denominator
\[
10,\ 100,\ 1000,\ \ldots
\]
which are all of the form $2^k5^k$. So after a fraction is written in lowest terms, its denominator can contain only $2$s and $5$s.

Conversely, if the denominator contains only $2$s and $5$s, we can multiply the numerator and denominator by enough extra $2$s or $5$s to make the denominator a power of $10$, and then the decimal stops. Any other prime factor can never be removed, so such a fraction recurs.
\end{enumerate}

\medskip
\noindent\textbf{Q16.}
\wsol
\begin{enumerate}[label=(\alph*),leftmargin=*]
\item
\[
\frac27=0.\dot{2}8571\dot{4},
\]
\[
\frac37=0.\dot{4}2857\dot{1},
\]
\[
\frac47=0.\dot{5}7142\dot{8},
\]
\[
\frac57=0.\dot{7}1428\dot{5},
\]
\[
\frac67=0.\dot{8}5714\dot{2}.
\]

\item Every seventh uses the same six digits in the same cyclic order:
\[
1\to4\to2\to8\to5\to7\to1.
\]
Each fraction simply starts at a different point in that cycle.

\item Let $x=0.\dot{2}8571\dot{4}=0.285714285714\ldots$.

The repeating block has $6$ digits, so multiply by $1\,000\,000$:
\[
1\,000\,000x=285\,714.285714285714\ldots
\]
Subtract $x$:
\[
999\,999x=285\,714.
\]
Therefore
\[
x=\frac{285\,714}{999\,999}.
\]
Dividing the numerator and denominator by $142\,857$ gives
\[
x=\frac27,
\]
which agrees with part (a).
\end{enumerate}

\medskip
\noindent\textbf{Q17.}
\wsol
\begin{enumerate}[label=(\alph*),leftmargin=*]
\item One valid example is $0.\dot{2}\dot{7}$.
Let $x=0.\dot{2}\dot{7}=0.272727\ldots$.
\[
100x=27.272727\ldots
\]
Subtract:
\[
99x=27,
\]
so
\[
x=\frac{27}{99}=\frac{3}{11}.
\]
So $0.\dot{2}\dot{7}$ is a recurring decimal with denominator $11$.

\item First,
\[
\frac56=0.8333\ldots=0.8\dot{3}.
\]
Check by turning it back: let $x=0.8\dot{3}=0.8333\ldots$.
\[
10x=8.3333\ldots,\qquad 100x=83.3333\ldots
\]
Subtract:
\[
100x-10x=83.3333\ldots-8.3333\ldots=75.
\]
Thus
\[
90x=75,
\]
so
\[
x=\frac{75}{90}=\frac56.
\]

\item Yes. Suppose the repeating block has $n$ digits and write the block as the whole number $N$. Shifting the decimal point by $n$ places gives
\[
10^n x=N+x,
\]
because the same repeating tail appears again. Therefore
\[
(10^n-1)x=N,
\]
so
\[
x=\frac{N}{10^n-1},
\]
which is a fraction.

If the decimal has some non-repeating digits at the front, first multiply by a power of $10$ to move past those digits. Then apply the same idea to the remaining recurring part. So the subtraction always works.

\item By part (c), every recurring decimal is equal to a fraction. But $\pi$ is not equal to a fraction; it is irrational. Therefore no recurring decimal can be exactly equal to $\pi$. Its decimal expansion goes on forever but never settles into a repeating block.
\end{enumerate}

\end{document}